\section{Automated Maturity Assessment Protocol}
\label{sec:assessment}

A persistent weakness in maturity model engineering is the reliance on manual, subjective Likert-scale questionnaires that suffer from respondent bias, high assessment labor costs, and rapid obsolescence in dynamic cloud environments \cite{proenca_2016, proenca_2016_rcis}. To address Research Question \textbf{RQ2}, this section formalizes an objective, programmatic maturity assessment protocol for the CN-OMM. Grounded in ISO/IEC 33001 process capability concepts \cite{iso_33001} and Lasrado et al.'s geometric distance metrics \cite{lasrado_2015}, our protocol translates the 15 sub-components of Table~\ref{tab:cn_omm_dimensions} into automated cluster audit queries and a non-compensatory mathematical scoring function.

\subsection{Declarative Cluster Auditing via Quantitative Capability Indicators}
\label{subsec:declarative_qci}

Rather than asking engineers whether their cluster is ``well instrumented,'' the CN-OMM assessment engine inspects the live Kubernetes API server, OpenTelemetry Collector configuration graphs, Prometheus/Mimir time-series metadata, and Chaos Mesh/LitmusChaos Custom Resource Definitions (CRDs). For each sub-component $d_{k.j}$ ($k \in \{1..5\}, j \in \{1..3\}$), the audit probe evaluates a normalized Quantitative Capability Indicator $\text{QCI}_{k.j} \in [0, 1]$, which maps onto the continuous capability scale $c_{k.j} \in [1.0, 5.0]$ via piecewise monotonic thresholds:
\begin{itemize}[leftmargin=*, nosep]
    \item \textbf{Telemetry \& Standardization Probes ($D_1$):} Queries the Kubernetes API for active Pods in target namespaces and cross-references OTLP collector receiver metrics (\texttt{otelcol\_receiver\_accepted\_spans}, \texttt{accepted\_metric\_points}, \texttt{accepted\_log\_records}) and eBPF DaemonSet status to compute $\text{QCI}_{1.1}, \text{QCI}_{1.2},$ and $\text{QCI}_{1.3}$.
    \item \textbf{Semantic Coupling Probes ($D_2$):} Samples structured log streams in Loki (`{namespace="target"} | json | trace_id != ""`) to compute the log-to-trace coupling ratio $\rho_{\text{log-trace}}$, queries Prometheus `/api/v1/query_exemplars` to verify active OpenMetrics Exemplar links ($\rho_{\text{exemplar}}$), and checks whether the trace analytics pipeline executes the exclusive self-time operator $T_{\text{self}}(s_i)$ ($\text{QCI}_{2.3}$).
    \item \textbf{Diagnostic \& Chaos Verification Probes ($D_3, D_4$):} Evaluates historical or continuous Chaos Engineering experiment objects (\texttt{PodChaos}, \texttt{NetworkChaos}, \texttt{StressChaos}, \texttt{ChaosEngine}) to verify the presence of declarative steady-state probes ($\text{QCI}_{4.2}$) and computes empirical distributions of $MTTD$, $MTTDiag$, $Acc$, and cascading false-positive rates ($\text{FPR}_{\text{cascade}}$).
    \item \textbf{Probe-Effect \& FinOps Probes ($D_5$):} Computes the steady-state marginal CPU and memory overhead $\Delta_{\text{CPU}}, \Delta_{\text{RAM}}$ of all observability components (\texttt{container\_cpu\_usage\_seconds\_total}) against the $\theta_{\max} = 15\%$ stability threshold ($\text{QCI}_{5.1}$) and audits storage index amplification ratios ($\text{QCI}_{5.3}$).
\end{itemize}

\subsection{Non-Compensatory Stage-Gate Scoring and Euclidean Progression Index}
\label{subsec:scoring_math}

Given the capability scores $c_{k.j} \in [1.0, 5.0]$ across the three sub-components of dimension $D_k$, the dimensional capability score $C(D_k)$ is given by:
\begin{equation}
    C(D_k) = \frac{1}{3} \sum_{j=1}^{3} c_{k.j}, \quad \forall k \in \{1, 2, 3, 4, 5\}
    \label{eq:dimension_score}
\end{equation}
yielding the five-dimensional capability profile vector $\mathbf{c}(\mathcal{S}) = [C(D_1), C(D_2), C(D_3), C(D_4), C(D_5)]^\top \in [1, 5]^5$.

In naive maturity models, the overall maturity level is calculated as the compensatory arithmetic mean $\bar{C} = \frac{1}{5}\sum_{k=1}^5 C(D_k)$. However, as established by configurational theory \cite{fiss_2011, lasrado_2015} and De Bruin et al.'s stage-gate principle \cite{debruin_2005}, an arithmetic mean masks catastrophic architectural bottlenecks: a cluster with high telemetry ingestion ($C(D_1) = 4.0$) but zero cross-pillar correlation ($C(D_2) = 2.0$) or runaway $30\%$ CPU overhead ($C(D_5) = 1.5$) cannot operate at Level 3 or Level 4 reliability. Therefore, CN-OMM defines the \textbf{Effective Stage-Gate Maturity Level} $\text{Level}_{\text{eff}}(\mathcal{S}) \in \{1, 2, 3, 4, 5\}$ using a strict non-compensatory minimum operator:
\begin{equation}
    \text{Level}_{\text{eff}}(\mathcal{S}) = \min_{k \in \{1, \dots, 5\}} \left\lfloor C(D_k) \right\rfloor
    \label{eq:stage_gate_level}
\end{equation}
Under Equation~(\ref{eq:stage_gate_level}), a system achieves maturity stage $L \in \{1..5\}$ if and only if every orthogonal dimension satisfies the necessary stage-gate predicate $C(D_k) \ge L$.

To provide granular visibility into continuous improvement between discrete stage-gates, we also formalize the normalized \textbf{Euclidean Maturity Progression Index} $\mathcal{I}(\mathcal{S}) \in [0, 1]$ relative to the ideal reference profile $\mathbf{c}^*(L_{\max}) = [5, 5, 5, 5, 5]^\top$ \cite{lasrado_2015}:
\begin{equation}
    \mathcal{I}(\mathcal{S}) = 1 - \frac{\left\| \mathbf{c}(\mathcal{S}) - \mathbf{c}^*(L_{\max}) \right\|_2}{\sqrt{5 \cdot (5 - 1)^2}} = 1 - \frac{\sqrt{\sum_{k=1}^{5} \left(5 - C(D_k)\right)^2}}{\sqrt{80}}
    \label{eq:euclidean_index}
\end{equation}
Together, the tuple $\langle \text{Level}_{\text{eff}}(\mathcal{S}), \mathcal{I}(\mathcal{S}), k_{\text{bottleneck}} \rangle$, where $k_{\text{bottleneck}} = \arg\min_k C(D_k)$, provides both an unambiguous certification of current stage-gate maturity and an actionable prescriptive roadmap identifying the exact limiting dimension that must be remediated next.
